% ============================================================
% CAPITOLO 37 - RIEPILOGO DELLA PARTE VII
% ============================================================
\chapter{Riepilogo della Parte VII}
\label{chap:riepilogo_vii}

\section{Tonicizzare o modulare?}
\label{sec:tonicizzare_o_modulare}

Il criterio decisivo è la cadenza di conferma (Sezione \ref{subsec:conferma_tonalita}).

\begin{center}
\small
\begin{tabular}{p{3.6cm}p{4.4cm}p{4.2cm}}
\toprule
 & \textbf{Tonicizzazione} & \textbf{Modulazione} \\
\midrule
Durata & uno o due accordi & stabile \\
Dopo l'accordo tonicizzato & si torna alla tonalità d'impianto & segue una cadenza V--I nella nuova tonalità \\
Segnale nel basso & alterazione isolata (sensibile di una dominante secondaria) & alterazioni che si ripetono e si stabilizzano \\
Dettaglio & Capitolo \ref{chap:tonicizzazione} & Capitoli \ref{chap:modulazione}--\ref{chap:progressioni_modulanti} \\
\bottomrule
\end{tabular}
\end{center}

\section{Tabella toni vicini e pivot}
\label{sec:tabella_toni_vicini}

\begin{center}
\small
\begin{tabular}{lp{6cm}p{3.4cm}}
\toprule
\textbf{Partenza} & \textbf{Toni vicini} & \textbf{Dettaglio} \\
\midrule
Maggiore & relativa minore, dominante e sua relativa, sottodominante e sua relativa & Sez. \ref{subsec:toni_vicini_maggiore} \\
Minore & relativa maggiore (III), IV, V, VI, VII & Sez. \ref{subsec:toni_vicini_minore} \\
\bottomrule
\end{tabular}
\end{center}

I pivot più efficaci sono descritti nelle Sezioni \ref{subsec:pivot_relative}, \ref{subsec:pivot_dominante} e \ref{subsec:pivot_sottodominante}.
